% TOP_PROBLEM: {"id": 30006851, "title": "Marked Length Spectrum Rigidity Conjecture", "category_id": 6, "category": {"id": 6, "name": "geometry", "display_name": "Geometry", "description": "Euclidean and non-Euclidean geometry, geometric structures.", "slug": "geometry", "order_index": 6, "created_at": "2026-07-31T15:26:25.671Z"}, "status": "open"}
% FIRST_POSTED: 2026-09-25
% ENTRY_KIND: statement_only
% !TeX program = lualatex
% Definition and sources only; this file is not an LLM solution attempt.
\documentclass[11pt]{article}
\usepackage[margin=25mm]{geometry}
\usepackage{fontspec}
\setmainfont{Latin Modern Roman}
\usepackage{amsmath,amssymb,mathtools}
\usepackage{unicode-math}
\setmathfont{Latin Modern Math}
\usepackage{xurl,hyperref}
\hypersetup{colorlinks=true,urlcolor=blue}
\setcounter{secnumdepth}{0}
\setlength{\parindent}{0pt}
\setlength{\parskip}{5pt}
\setlength{\emergencystretch}{3em}
\begin{document}
\section*{\texorpdfstring{Marked Length Spectrum Rigidity Conjecture}{Marked Length Spectrum Rigidity Conjecture}}\label{30006851}
\subsection{Definitions and mathematical statement}
Let $M$ be a closed manifold and let $g,h$ be negatively curved Riemannian metrics. Each nontrivial free homotopy class $[\gamma]$ of loops has a unique closed geodesic representative for each metric. Its length defines the \emph{marked length spectrum} $\ell_g([\gamma])$. The conjecture asserts
\[
 \ell_g([\gamma])=\ell_h([\gamma])\text{ for every }[\gamma]
 \quad\Longrightarrow\quad
 \exists f:M\to M\text{ isotopic to }\mathrm{id},\quad f^*h=g.
\]
The marking identifies classes in the same manifold, not just an unordered collection of numerical lengths. The isotopy requirement is part of the supplied target and should not be dropped in favour of an arbitrary isometry. Negative curvature refers to sectional curvature in every tangent two-plane.
\par\smallskip\textit{Formulation and concept references:} \href{https://aimath.org/pastworkshops/geodesicsproblems.pdf}{[S1]}.

\subsection{Short English statement}
Do the lengths of closed geodesics, labelled by their loop classes, determine a negatively curved metric up to an isometry isotopic to the identity?
\subsection{Sources}
\begin{itemize}
\item[S1]AIM workshop problem list, Open problems and questions about geodesics.
\url{https://aimath.org/pastworkshops/geodesicsproblems.pdf}.
\textit{Formulation source.}
\end{itemize}
\end{document}
